\scanpage{244}
\mspage{118}
% archive sheet number: 243

% slanted
\margin{J'ai changé la numérotation, pour tenir compte d'un Loc (7) plus bas, qui est encore un axiome de stabilité.}

\uline{Loc \deleted{\unsure{7}}8} (Axiome de connexité) $f\in W\Rightarrow\pi_0(f)$ bijectif

\uline{NB} Tous ces axiomes, sauf le dernier (et un \unsure{chouïa|chouia} de l'avant-dernier) \unsure{sont} des axiomes de stabilité --- ils disent tous que $W$ est ``grand'' : \deleted{\unsure{S}}il contient certaines flèches, et s'il contient telles flèches, il contient telles autres.

Loc \deleted{\unsure{7}}8, et dans une certaine mesure Loc \deleted{\unsure{6}}7 bis, sont des axiomes ``limitatifs'' sur la taille de $W$ : Loc 7 nous dit que $W\subset W_0$ ($=$ ens des $f\in\mathrm{Fl}\,\mathrm{Cat}$ telles que $\pi_0(f)$ bij.), Loc \deleted{\unsure{6}}7 bis nous dit l'existence d'un cardinal tel que $W$ soit égal à la \inserted{$\unsure{L|k}$-}\unsure{enveloppe} de $W^\pi$ (donc $W\subset\unsure{\mathrm{l}\text{-}\mathrm{env}|\mathrm{L}\text{-}\mathrm{env}}(W^\pi)$).

J'ai \unsure{essayé|envisagé} de ranger les axiomes par ordre d'importance pratique. J'ai fini par \unsure{rejeter|regarder} l'axiome de connexité à la fin, car dans les développements théoriques sur les localiseurs généraux, il n'intervient pratiquement jamais.

Mais je rappelle (\uline{XII} p.~\unsure{10|20}, Th.~2) que si $W$ satisfait $L(1,2,3,\unsure{6|5})$, et si $L(\deleted{\unsure{7}}8)$ n'est \uline{pas} satisfait, \deleted{\unsure{il suffit à 5 b}} alors $W=\mathrm{Fl}(\mathrm{Cat})$ ; \deleted{\unsure{donc}} en présence de $L(1,2,3,5)$, $L(\deleted{\unsure{7}}8)$ \uline{équivaut} à $W\ne\mathrm{Fl}(\mathrm{Cat})$ tout entier.
